\section{Etapa 11 --- La placa nueva (rama \texttt{cambios-hardware})}
\label{sec:et11}
\periodo{2 --- 12 de agosto de 2026 \quad$\cdot$\quad etapa en curso}

\dato{La bit\'acora diaria termina el 1\textsuperscript{o} de agosto. Esta
secci\'on est\'a reconstruida directamente de los \emph{commits} y de la
documentaci\'on de construcci\'on del firmware, no de la bit\'acora.}

Se dise\~n\'o una placa propia ---hasta ac\'a todo hab\'ia corrido sobre
\emph{kits} de desarrollo y una \emph{carrier} intermedia
(CY8CKIT-059 $+$ ESP32 DevKitC)--- y eso oblig\'o a rehacer el TopDesign del
PSoC y buena parte del enlace entre procesadores. Todo el trabajo vive en la
rama \texttt{cambios-hardware}, deliberadamente separada de \texttt{main}.

\subsection{Cambios en el PSoC}

\begin{center}\small
\begin{tabularx}{\textwidth}{@{}l X@{}}
\toprule
\textbf{Cambio} & \textbf{Motivo y consecuencia} \\
\midrule
\textbf{IDAC en lugar de VDAC} en las cuatro referencias de calibraci\'on &
Se pasa a fijar las referencias por \emph{corriente} en vez de por tensi\'on,
con las globales de placa $R = \SI{30}{\kilo\ohm}$,
$V_{\mathrm{ref}} = \SI{2.048}{\volt}$,
$I_{fs} = \SI{31.875}{\micro\ampere}$ y helpers de conversi\'on
c\'odigo$\leftrightarrow$tensi\'on. \textbf{El LSB queda en
\SI{3.75}{\milli\volt}}, que es la nueva resoluci\'on del actuador del lazo PI
de la Etapa~6 \\
\addlinespace[3pt]
\textbf{La UART queda solo como RX} &
El TX desaparece del TopDesign nuevo \\
\addlinespace[3pt]
\textbf{Subida PSoC$\to$ESP por I\textsuperscript{2}C} (\texttt{psoc\_link\_*}, PSoC maestro) &
\emph{La UART a 115\,200 era el cuello de botella del \emph{stream} de
muestras.} El protocolo de \emph{frames} no cambia: cambia el medio \\
\addlinespace[3pt]
\textbf{\texttt{PGAout}}: nueva etapa de salida del \emph{pipeline} GEO &
Configurable por el comando \texttt{0xA8}, igual que \texttt{PGAgain} y \texttt{PGAvdac} \\
\bottomrule
\end{tabularx}
\end{center}

\seleccionado{\textbf{El principio de la Etapa~4 sobrevive intacto
al cambio de placa}: el medio del enlace cambi\'o por tercera vez
(SPI $\to$ UART $\to$ I\textsuperscript{2}C) sin tocar el protocolo. Que el
formato de \emph{frames} sea independiente del transporte es lo que hizo
barata cada migraci\'on.}

\subsubsection{Trampas del TopDesign nuevo, registradas}

\pendiente{Estas cost\'o encontrarlas y conviene no volver a tropezar:
\begin{itemize}[itemsep=0.15em,topsep=0.2em]
\item El componente se llama \texttt{I2C\_1}, \textbf{no} \texttt{I2C}.
\item \textbf{El TopDesign nuevo no tiene pin LED}: hubo que sacar el
\emph{include} directo de \texttt{LED.h} para que los \emph{guards}
\texttt{CY\_PINS\_LED\_H} hicieran efecto.
\item Los pines de \texttt{SPIp} quedaron \textbf{repartidos en varios
puertos}, as\'i que el \emph{fitter} dej\'o de emitir los macros agregados; para
el CS de la SD se usan los del pin~0.
\end{itemize}
Y sigue vigente la regla del proyecto: \textbf{nunca editar a mano}
\texttt{.cyprj}, \texttt{.cydwr}, \texttt{.cysch} ni \texttt{.cyfit} ---solo
\texttt{.c} y \texttt{.h}. Romper esto ya rompi\'o un \emph{build} una vez.}

\dato{Estado del \emph{build} tras adaptar el firmware al TopDesign nuevo:
\textbf{compila}. Flash \SI{60400}{\byte}, SRAM \SI{51928}{\byte}.}

\subsection{Cambios en el ESP32}

\textbf{El enlace pasa a ser asim\'etrico.} ESP\,$\to$\,PSoC sigue por UART
(porque el PSoC conserva el RX); PSoC\,$\to$\,ESP pasa a I\textsuperscript{2}C
con el PSoC de maestro y el ESP de esclavo en \texttt{PSOC\_I2C\_ADDR}
(GPIO\,21/22). Los bytes entran a un \emph{ring buffer} y \texttt{poll()} los
consume con el mismo \emph{parser} de \emph{frames} de siempre.

\textbf{Sincronizaci\'on: se elimina el cable.} Desaparecen
\texttt{SYNC\_IN\_PIN} y la ISR \texttt{onSyncEdge}. Ya no hay sincronizaci\'on
por cable desde el maestro: \textbf{el arranque llega \'unicamente por ESP-NOW},
y es el propio esclavo el que genera el flanco hacia el PSoC por
\texttt{SYNC\_TO\_PSOC\_PIN}, el \'unico pin de \emph{sync} que queda.

\seleccionado{Esto simplifica el despliegue en campo ---un cable menos por
nodo--- \textbf{sin violar el principio de instante cero por flanco f\'isico}:
el flanco sigue existiendo, solo que ahora lo genera el esclavo localmente al
recibir el disparo por radio. PRESTART/HOT\_WAIT garantiza que todos est\'en
listos; la dispersi\'on temporal entre nodos queda gobernada por la entrega y el
despacho de \texttt{CMD\_START} m\'as el camino GPIO$\to$PSoC.}

\pendiente{Consecuencia directa: \textbf{el presupuesto de \emph{jitter} hay
que rehacerlo} para esta topolog\'ia. Con \emph{sync} por cable, una parte del
error estaba acotada por hardware; ahora todo el error de sincronizaci\'on pasa
por el enlace de radio. Es la verificaci\'on m\'as urgente de esta etapa.}

El firmware incluye un ensayo reproducible para cerrarla:
\texttt{SCOPE\_MULTI} ejecuta 128 arranques, separados \SI{1}{\second},
repitiendo el PRESTART completo. Para interpretar la distribuci\'on medida se
usa $\Delta\phi=2\pi f\Delta t$: \SI{400}{\micro\second} equivalen a
\SI{7.2}{\degree} a \SI{50}{\hertz} y \SI{14.4}{\degree} a
\SI{100}{\hertz}; un l\'imite de \SI{5}{\degree} exige
\SI{278}{\micro\second} y \SI{139}{\micro\second}, respectivamente.

\begin{figure}[!ht]
\centering
\includegraphics[width=0.86\textwidth]{sincronizacion_fase.png}
\caption{Presupuesto de sincronizaci\'on expresado en la variable que importa
para MASW: fase. La observaci\'on menor a \SI{400}{\micro\second} es preliminar
y no sustituye la distribuci\'on de \texttt{SCOPE\_MULTI}.}
\label{fig:sync-fase-hist}
\end{figure}

\textbf{Interfaz local en el nodo} (OLED por SPI $+$ botones). Cada esclavo
puede operarse sin maestro y sin tel\'efono. Junto con los comandos USB de la
Etapa~6, completa tres v\'ias independientes de operaci\'on:
interfaz local, USB y web.

\textbf{\texttt{PGAout} en toda la cadena de control}: comando \texttt{0xA8},
configuraci\'on dirigida por ESP-NOW, comando USB \texttt{pgaout N}, fila propia
en el panel web ---visible solo cuando el PSoC reporta
\texttt{hw\_class=GEO}--- con persistencia, y exportaci\'on como
\texttt{pgaout\_code}/\texttt{pgaout\_gain} en los metadatos.

\seleccionado{Que un par\'ametro nuevo aparezca simult\'aneamente en el firmware,
en los tres frentes de operaci\'on \emph{y} en los metadatos exportados es la
se\~nal de que la arquitectura de configuraci\'on madur\'o: agregar una etapa ya
no es un parche, es rellenar un patr\'on.}

\subsection{PCBs}

El repositorio \texttt{PCBs} contiene el proyecto en \textbf{KiCad} y en
\textbf{JitX} (m\'as un \texttt{JitX\_para\_port}), con \emph{commits} hasta el
9 de agosto.

\pendiente{No hay bit\'acora de esta etapa y la decisi\'on KiCad frente a JitX
no est\'a registrada en ning\'un lado. \textbf{¿Cu\'al es el proyecto vigente y por
qu\'e se explor\'o el otro?}}

\subsection{Balance de la etapa}

\begin{center}\small
\begin{tabularx}{\textwidth}{@{}l X@{}}
\toprule
\textbf{Se qued\'o} & Placa propia; IDAC como referencia de calibraci\'on (LSB \SI{3.75}{\milli\volt}); enlace asim\'etrico UART bajada / I\textsuperscript{2}C subida; etapa \texttt{PGAout}; interfaz local OLED $+$ botones; \emph{sync} \'unicamente por ESP-NOW con flanco generado localmente \\
\addlinespace[3pt]
\textbf{Se descart\'o} & VDAC como referencia; TX de la UART del PSoC; \emph{sync} por cable entre maestro y esclavo (\texttt{SYNC\_IN\_PIN} y su ISR); pin de LED en el TopDesign \\
\addlinespace[3pt]
\textbf{Qued\'o abierto} & Rehacer el presupuesto de \emph{jitter} sin cable de \emph{sync}; recalibrar el lazo PI con la nueva resoluci\'on del IDAC; validar la placa contra la matriz E1--E18 de la Etapa~8; decidir KiCad frente a JitX \\
\bottomrule
\end{tabularx}
\end{center}
